% ============================================================================
% Презентация на курсовия проект по „Сигурни протоколи и архитектури за комуникации“.
% Компилация (от docs/slides/):  latexmk Slides.tex
% Числата идват от ../generated/, същите файлове, които ползва документът. Листингите
% се четат направо от изходния код. Нищо тук не е преписано на ръка от записите.
% ============================================================================
\documentclass[aspectratio=169, 11pt]{beamer}

\usepackage{cmap}
\usepackage[T2A]{fontenc}
\usepackage[utf8]{inputenc}
\usepackage[bulgarian]{babel}
\usepackage{paratype}
\renewcommand{\familydefault}{\sfdefault}
\usepackage{booktabs}
\usepackage{tikz}
\usetikzlibrary{arrows.meta,positioning,calc}
\usepackage{listings}

% Generated by bench/analyse.py. Do not edit by hand.
\newcommand{\winTlsDiffA}{6813{,}5}
\newcommand{\winTlsLowA}{5078{,}0}
\newcommand{\winTlsHighA}{8533{,}0}
\newcommand{\winTlsRelA}{33{,}4}
\newcommand{\winClearDiffA}{4332{,}5}
\newcommand{\winClearLowA}{2950{,}0}
\newcommand{\winClearHighA}{6150{,}5}
\newcommand{\winClearRelA}{35{,}4}
\newcommand{\winTlsDiffS}{28{,}5}
\newcommand{\winTlsLowS}{5{,}0}
\newcommand{\winTlsHighS}{53{,}0}
\newcommand{\winTlsRelS}{0{,}9}
\newcommand{\winClearDiffS}{-2{,}0}
\newcommand{\winClearLowS}{-4{,}0}
\newcommand{\winClearHighS}{1{,}0}
\newcommand{\winClearRelS}{-1{,}6}
\newcommand{\winTlsDetectMedian}{3{,}137}
\newcommand{\winTlsDetectDecideMedian}{23{,}869}
\newcommand{\winTlsDetectStartMedian}{23{,}048}
\newcommand{\winTlsDetectStartNinety}{81{,}672}
\newcommand{\winTlsDedicatedMedian}{3{,}108}
\newcommand{\winTlsDedicatedDecideMedian}{14{,}893}
\newcommand{\winTlsDedicatedStartMedian}{14{,}893}
\newcommand{\winTlsDedicatedStartNinety}{73{,}974}
\newcommand{\winClearDetectMedian}{0{,}126}
\newcommand{\winClearDetectDecideMedian}{16{,}471}
\newcommand{\winClearDetectStartMedian}{14{,}305}
\newcommand{\winClearDetectStartNinety}{64{,}988}
\newcommand{\winClearDedicatedMedian}{0{,}128}
\newcommand{\winClearDedicatedDecideMedian}{9{,}747}
\newcommand{\winClearDedicatedStartMedian}{9{,}746}
\newcommand{\winClearDedicatedStartNinety}{69{,}500}
\newcommand{\winReps}{5}
\newcommand{\winConns}{400}
\newcommand{\winTotal}{2000}
\newcommand{\winRate}{100}
\newcommand{\winDuration}{273}
\newcommand{\winCommit}{c08784d}
\newcommand{\winSeed}{20261103}
\newcommand{\winPlatform}{Windows-11-10.0.26200-SP0}
\newcommand{\winFirstReadMedian}{517}
\newcommand{\winServed}{8000}
\newcommand{\winMeasured}{8000}

% Generated by bench/analyse.py. Do not edit by hand.
\newcommand{\linTlsDiffA}{16{,}5}
\newcommand{\linTlsLowA}{1{,}0}
\newcommand{\linTlsHighA}{31{,}5}
\newcommand{\linTlsRelA}{0{,}7}
\newcommand{\linClearDiffA}{10{,}0}
\newcommand{\linClearLowA}{6{,}0}
\newcommand{\linClearHighA}{12{,}5}
\newcommand{\linClearRelA}{4{,}2}
\newcommand{\linTlsDiffS}{6{,}0}
\newcommand{\linTlsLowS}{-8{,}0}
\newcommand{\linTlsHighS}{20{,}5}
\newcommand{\linTlsRelS}{0{,}3}
\newcommand{\linClearDiffS}{5{,}0}
\newcommand{\linClearLowS}{2{,}0}
\newcommand{\linClearHighS}{8{,}0}
\newcommand{\linClearRelS}{6{,}3}
\newcommand{\linTlsDetectMedian}{2{,}169}
\newcommand{\linTlsDetectDecideMedian}{0{,}336}
\newcommand{\linTlsDetectStartMedian}{0{,}158}
\newcommand{\linTlsDetectStartNinety}{4{,}543}
\newcommand{\linTlsDedicatedMedian}{2{,}163}
\newcommand{\linTlsDedicatedDecideMedian}{0{,}165}
\newcommand{\linTlsDedicatedStartMedian}{0{,}165}
\newcommand{\linTlsDedicatedStartNinety}{0{,}358}
\newcommand{\linClearDetectMedian}{0{,}084}
\newcommand{\linClearDetectDecideMedian}{0{,}172}
\newcommand{\linClearDetectStartMedian}{0{,}159}
\newcommand{\linClearDetectStartNinety}{0{,}295}
\newcommand{\linClearDedicatedMedian}{0{,}079}
\newcommand{\linClearDedicatedDecideMedian}{0{,}153}
\newcommand{\linClearDedicatedStartMedian}{0{,}152}
\newcommand{\linClearDedicatedStartNinety}{2{,}863}
\newcommand{\linReps}{5}
\newcommand{\linConns}{400}
\newcommand{\linTotal}{2000}
\newcommand{\linRate}{100}
\newcommand{\linDuration}{101}
\newcommand{\linCommit}{c08784d}
\newcommand{\linSeed}{20261103}
\newcommand{\linPlatform}{Linux-6.6.114.1-microsoft-standard-WSL2-x86\_64-with-glibc2.39}
\newcommand{\linFirstReadMedian}{517}
\newcommand{\linServed}{8000}
\newcommand{\linMeasured}{8000}

% Generated by bench/analyse.py. Do not edit by hand.
\newcommand{\winfirstTlsDiffA}{280{,}0}
\newcommand{\winfirstTlsLowA}{81{,}0}
\newcommand{\winfirstTlsHighA}{496{,}0}
\newcommand{\winfirstTlsRelA}{6{,}3}
\newcommand{\winfirstClearDiffA}{83{,}5}
\newcommand{\winfirstClearLowA}{7{,}0}
\newcommand{\winfirstClearHighA}{162{,}5}
\newcommand{\winfirstClearRelA}{13{,}8}
\newcommand{\winfirstTlsDetectMedian}{4{,}726}
\newcommand{\winfirstTlsDetectDecideMedian}{1{,}323}
\newcommand{\winfirstTlsDedicatedMedian}{4{,}446}
\newcommand{\winfirstTlsDedicatedDecideMedian}{0{,}749}
\newcommand{\winfirstClearDetectMedian}{0{,}688}
\newcommand{\winfirstClearDetectDecideMedian}{0{,}617}
\newcommand{\winfirstClearDedicatedMedian}{0{,}604}
\newcommand{\winfirstClearDedicatedDecideMedian}{0{,}396}
\newcommand{\winfirstReps}{5}
\newcommand{\winfirstConns}{400}
\newcommand{\winfirstTotal}{2000}
\newcommand{\winfirstRate}{100}
\newcommand{\winfirstDuration}{98}
\newcommand{\winfirstCommit}{69b599e}
\newcommand{\winfirstSeed}{20261103}
\newcommand{\winfirstPlatform}{Windows-11-10.0.26200-SP0}
\newcommand{\winfirstFirstReadMedian}{517}
\newcommand{\winfirstServed}{8000}
\newcommand{\winfirstMeasured}{8000}

% Generated by bench/evidence2tex.py. Do not edit by hand.
\newcommand{\microWinDedicated}{12{,}8}
\newcommand{\microWinDetecting}{224{,}3}
\newcommand{\microWinDiff}{211{,}5}
\newcommand{\microWinRounds}{15}
\newcommand{\microWinIterations}{200000}
\newcommand{\microLinDedicated}{36{,}0}
\newcommand{\microLinDetecting}{199{,}6}
\newcommand{\microLinDiff}{163{,}6}
\newcommand{\microLinRounds}{15}
\newcommand{\microLinIterations}{200000}
\newcommand{\checksClassify}{538}
\newcommand{\checksReplay}{12}
\newcommand{\checksHttp}{22}
\newcommand{\checksTls}{47}
\newcommand{\checksIntegration}{35}
\newcommand{\checksUnitTotal}{619}

% Generated by bench/plot2tex.py. Do not edit by hand.
\def\plotwinTlsDiff{28.5}
\def\plotwinTlsLow{5.0}
\def\plotwinTlsHigh{53.0}
\def\plotwinClearDiff{-2.0}
\def\plotwinClearLow{-4.0}
\def\plotwinClearHigh{1.0}
\def\plotlinTlsDiff{6.0}
\def\plotlinTlsLow{-8.0}
\def\plotlinTlsHigh{20.5}
\def\plotlinClearDiff{5.0}
\def\plotlinClearLow{2.0}
\def\plotlinClearHigh{8.0}


% --- Цветове: един акцент, и по един цвят за всеки изход на разпознаването -----------
\definecolor{ink}{RGB}{28,32,38}
\definecolor{muted}{RGB}{110,116,124}
\definecolor{rule}{RGB}{214,218,224}
\definecolor{accent}{RGB}{0,74,128}
\definecolor{tls}{RGB}{0,128,96}
\definecolor{clear}{RGB}{196,110,0}
\definecolor{unknown}{RGB}{178,34,34}
\colorlet{tlsbg}{tls!10}
\colorlet{clearbg}{clear!12}
\colorlet{unknownbg}{unknown!9}

% --- Оформление -----------------------------------------------------------------
\setbeamertemplate{navigation symbols}{}
\setbeamercolor{normal text}{fg=ink}
\setbeamercolor{structure}{fg=accent}
\setbeamercolor{alerted text}{fg=unknown}
\setbeamerfont{frametitle}{size=\Large, series=\bfseries}
\setbeamertemplate{frametitle}{%
  \vspace{0.55em}\insertframetitle\par
  \vspace{0.15em}{\color{accent}\rule{2.2em}{2pt}}\vspace{0.1em}}
\setbeamertemplate{footline}{%
  \hspace{1em}{\color{muted}\scriptsize octet: TLS и HTTP на един порт}%
  \hfill{\color{muted}\scriptsize\insertframenumber\,/\,\inserttotalframenumber}\hspace{1em}\vspace{0.7em}}
\setbeamertemplate{itemize items}{\color{accent}\small$\blacktriangleright$}
\setbeamertemplate{itemize subitem}{\color{muted}\small$\cdot$}
\setbeamersize{text margin left=1.4em, text margin right=1.4em}

\newcommand{\code}[1]{{\ttfamily #1}}
\newcommand{\hex}[1]{{\ttfamily\bfseries #1}}
\newcommand{\src}[1]{{\color{muted}\scriptsize #1}}
% One row of the effects chart: y, label, interval low, interval high, colour.
\newcommand{\effect}[5]{%
  \node[anchor=east, font=\scriptsize] at (\xmin-1,#1) {#2};
  \draw[#5, line width=2.4pt] (#3,#1) -- (#4,#1);}

\lstset{
  language=C++, basicstyle=\small\ttfamily, columns=fullflexible, keepspaces=true,
  keywordstyle=\color{accent}\bfseries, commentstyle=\color{muted}, stringstyle=\color{clear},
  numbers=left, numberstyle=\tiny\color{muted}, numbersep=6pt, xleftmargin=2em,
  frame=l, rulecolor=\color{rule}, framerule=1pt, tabsize=2, showstringspaces=false,
  breaklines=true, breakindent=2em,
  inputencoding=utf8, morekeywords={constexpr, noexcept, span},
}

\tikzset{
  box/.style={draw=rule, rounded corners=3pt, align=center, font=\small, inner sep=5pt,
    minimum height=0.95cm, fill=white, line width=0.8pt},
  tlsbox/.style={box, draw=tls, fill=tlsbg},
  clearbox/.style={box, draw=clear, fill=clearbg},
  unkbox/.style={box, draw=unknown, fill=unknownbg},
  arr/.style={-{Latex[length=2.2mm]}, line width=0.9pt, color=ink!70},
  byte/.style={draw=ink!60, minimum width=1.25cm, minimum height=1cm, font=\ttfamily\bfseries\large},
}

\title{Разпознаване на TLS и некриптиран HTTP по първия октет при един слушащ сокет}
\subtitle{Механизъм и анализ на сигурността}
\author{Алекс Цветанов}
\date{2026}

\begin{document}

% ---------------------------------------------------------------------------------
\begin{frame}[plain]
\begin{tikzpicture}[remember picture, overlay]
  \fill[accent] (current page.north west) rectangle ([yshift=-0.5cm]current page.north east);
\end{tikzpicture}
\vspace{1.2em}
\begin{columns}[T]
\column{0.13\textwidth}
\includegraphics[width=\linewidth]{../images/tu-logo.jpeg}
\column{0.83\textwidth}
{\small\color{muted} Технически университет, София \,·\, Факултет по телекомуникации\par
Курсов проект по „Сигурни протоколи и архитектури за комуникации“\par}
\end{columns}
\vspace{2.2em}
{\LARGE\bfseries Разпознаване на TLS и некриптиран HTTP\\[0.2em] по първия октет при един слушащ сокет\par}
\vspace{0.6em}
{\large\color{accent} Механизъм и анализ на сигурността\par}
\vspace{2.4em}
{\normalsize Алекс Цветанов, група 263\par}
{\small\color{muted} Ръководител: проф. Георги Илиев\par}
\vfill
{\small\color{muted}\ttfamily github.com/Alex-Tsvetanov/octet\hfill София, 2026}
\end{frame}

% ---------------------------------------------------------------------------------
\begin{frame}{Два порта или един}
\begin{columns}[T]
\column{0.48\textwidth}
\textbf{Обичайно}\par\medskip
\begin{tikzpicture}[node distance=0.5cm]
\node[clearbox, minimum width=4.6cm] (p80) {порт 80: некриптиран HTTP};
\node[tlsbox, minimum width=4.6cm, below=of p80] (p443) {порт 443: HTTP над TLS};
\end{tikzpicture}\par\medskip
{\small Всеки порт знае предварително какво ще пристигне.\par}
\column{0.48\textwidth}
\textbf{В този проект}\par\medskip
\begin{tikzpicture}
\node[box, minimum width=4.6cm, line width=1.2pt, draw=accent] (one) {един слушащ сокет};
\node[below=0.25cm of one, font=\small, align=center] {за всяка връзка сървърът\\сам решава: TLS или HTTP?};
\end{tikzpicture}
\end{columns}
\vspace{1.4em}
\begin{itemize}
\item Решението се взема \alert{преди} TLS ръкостискането, когато нищо още не е защитено.
\item Въпросите: стига ли един октет, губят ли се данни, какво струва, какво се отваря за нападателя.
\end{itemize}
\end{frame}

% ---------------------------------------------------------------------------------
\begin{frame}{Защо един октет е достатъчен}
\begin{columns}[T]
\column{0.5\textwidth}
\textbf{TLS} \src{RFC 8446, раздел 5.1}\par
Клиентът говори пръв. Първият запис е ClientHello, ContentType \hex{22 = 0x16}.\par\medskip
\textbf{HTTP} \src{RFC 9110, раздел 9.1}\par
Всеки метод е с главни латински букви: \code{GET}, \code{POST}, \dots\ Префиксът на
HTTP/2 започва с \code{PRI}. \src{RFC 9113, раздел 3.4}\par\medskip
Двете множества не се пресичат. Тестът минава през всичките 256 стойности
(\checksClassify{} проверки).
\column{0.46\textwidth}
\centering
\begin{tikzpicture}[x=0.3cm, y=0.3cm]
\foreach \h [count=\r from 0] in {0,1,2,3,4,5,6,7,8,9,A,B,C,D,E,F}{
  \node[font=\tiny\ttfamily, color=muted, anchor=east] at (-0.5, -\r) {\h0};
  \node[font=\tiny\ttfamily, color=muted] at (\r, 1) {\h};}
\foreach \v in {0,...,255}{
  \pgfmathtruncatemacro{\row}{div(\v,16)}
  \pgfmathtruncatemacro{\col}{mod(\v,16)}
  \ifnum\v=22 \def\f{tls}\else
    \ifnum\v>64 \ifnum\v<91 \def\f{clear}\else\def\f{unknownbg}\fi\else\def\f{unknownbg}\fi
  \fi
  \fill[\f] (\col-0.45, -\row-0.45) rectangle (\col+0.45, -\row+0.45);
}
\end{tikzpicture}\par
{\scriptsize
\textcolor{tls}{\rule{0.7em}{0.7em}} \hex{0x16}: TLS\quad
\textcolor{clear}{\rule{0.7em}{0.7em}} \hex{A}\dots\hex{Z}: HTTP\quad
\textcolor{unknownbg}{\rule{0.7em}{0.7em}} останалите: затваряне\par}
\end{columns}
\end{frame}

% ---------------------------------------------------------------------------------
\begin{frame}[fragile]{Разпознаването в кода}
\begin{columns}[T]
\column{0.6\textwidth}
\vspace{-0.6em}
\lstinputlisting[firstline=36, lastline=52, firstnumber=36, gobble=1, basicstyle=\scriptsize\ttfamily]{../../include/octet/classify.hpp}
\column{0.36\textwidth}
\src{include/octet/classify.hpp}\par\medskip
\begin{itemize}
\item Гледа само първия октет.
\item \code{constexpr} и \code{noexcept}: без памет, без изключения, проверимо при компилиране.
\item Празен вход дава „непознато“, а не предположение.
\end{itemize}
\end{columns}
\end{frame}

% ---------------------------------------------------------------------------------
\begin{frame}{Пътят на една връзка}
\centering
\begin{tikzpicture}[node distance=0.55cm and 0.75cm]
\node[box] (acc) {\code{accept()}};
\node[box, right=of acc] (read) {четене до\\2048 октета};
\node[box, right=of read, draw=accent, line width=1.2pt] (cls) {първият\\октет};
\node[tlsbox, right=1.3cm of cls, yshift=1.35cm] (tls) {TLS ръкостискане\\(OpenSSL)};
\node[clearbox, right=1.3cm of cls] (http) {анализатор на HTTP};
\node[unkbox, right=1.3cm of cls, yshift=-1.35cm] (unk) {затваряне без отговор};
\node[tlsbox, right=0.6cm of tls] (http2) {HTTP над TLS};
\draw[arr] (acc) -- (read);
\draw[arr] (read) -- (cls);
\draw[arr] (cls.east) -- node[above, sloped, font=\scriptsize\ttfamily] {0x16} (tls.west);
\draw[arr] (cls.east) -- node[above, font=\scriptsize\ttfamily] {A..Z} (http.west);
\draw[arr] (cls.east) -- node[below, sloped, font=\scriptsize] {друго} (unk.west);
\draw[arr] (tls) -- (http2);
\draw[line width=2pt, accent!60] ([yshift=0.45cm]acc.west |- tls.north) -- ([yshift=0.45cm]tls.north east)
  node[pos=0.5, above, font=\scriptsize, color=accent] {един общ краен срок: от \code{accept()} до края на ръкостискането};
\end{tikzpicture}
\vspace{1.2em}
\begin{itemize}
\item Нишка за всяка връзка, с горна граница на едновременните връзки.
\item Непознатото се затваря мълчаливо: отговорът би дал информация на сканиращия.
\end{itemize}
\end{frame}

% ---------------------------------------------------------------------------------
\begin{frame}{Прочетените октети не се губят}
\begin{columns}[T]
\column{0.47\textwidth}
\textbf{\code{MSG\_PEEK}}\par
Октетите остават в ядрото, следващото четене ги получава отново.
Цена: второ системно извикване за същите данни.
\column{0.47\textwidth}
\textbf{Връщане (\code{ReplayReader})} \src{избраният начин}\par
Октетите се четат наистина и се подават първи на TLS слоя, после се чете от сокета.
\end{columns}
\vspace{1em}
\begin{itemize}
\item OpenSSL работи върху два буфера в паметта (\code{BIO\_s\_mem}), а не върху сокета.
\item \code{BIO\_set\_mem\_eof\_return(-1)}: празен буфер означава „опитай пак“, не „край на потока“.
\item Всички \winTotal{} измерени TLS връзки с разпознаване: първото четене донася целия
  ClientHello, \winFirstReadMedian{} октета. Няма допълнително системно извикване.
\end{itemize}
\end{frame}

% ---------------------------------------------------------------------------------
\begin{frame}[fragile]{Един абсолютен краен срок}
\begin{columns}[T]
\column{0.42\textwidth}
Относителен срок, който започва отначало при всяко четене, не спира клиент, който праща
по един октет точно преди изтичането. Той държи връзката отворена неограничено.\par\medskip
Срокът в \code{octet} е \alert{момент във времето}, зададен при приемането. Той покрива
разпознаването \textbf{и} цялото ръкостискане.\par\medskip
{\small Клиент, който праща \hex{0x16} и спира, получава същия срок като клиент, който не праща нищо.\par}
\column{0.56\textwidth}
\lstinputlisting[firstline=27, lastline=37, firstnumber=27, gobble=2, basicstyle=\scriptsize\ttfamily]{../../include/octet/deadline.hpp}
\src{include/octet/deadline.hpp}
\end{columns}
\end{frame}

% ---------------------------------------------------------------------------------
\begin{frame}{Настройки на TLS}
\small
\renewcommand{\arraystretch}{1.15}
\begin{tabular}{@{}p{0.3\textwidth}p{0.66\textwidth}@{}}
\toprule
\textbf{Настройка} & \textbf{Причина} \\
\midrule
минимум TLS 1.2, по избор 1.3 & по-старите версии не се договарят \src{RFC 9325, 3.1.1} \\
без предоговаряне & атаки върху предоговарянето в TLS 1.2 \src{RFC 5746; RFC 7457, 2.10} \\
без компресия & отваря атаки като CRIME \src{RFC 7457, 2.6; RFC 9325, 3.3} \\
сървърът избира шифъра & не клиентът, от общите шифри \\
строг ALPN & без общ протокол: фатално предупреждение 120 \src{RFC 9325, 3.8} \\
ключът спрямо сертификата & грешката спира сървъра при стартиране, не при първия клиент \\
\bottomrule
\end{tabular}\par\medskip
{\small Тестов удостоверяващ орган, ECDSA P-256: CA и сертификат на сървъра за \code{localhost}.
Клиентът проверява веригата и името, без да се изключва нищо.\par}
\end{frame}

% ---------------------------------------------------------------------------------
\begin{frame}{Отказ по ALPN, видян в мрежата}
Клиент, чийто списък с протоколи няма \code{http/1.1}, получава преди ServerHello, некриптирано:
\vspace{1em}
\begin{center}
\begin{tikzpicture}[node distance=0pt]
\node[byte, fill=unknownbg] (b0) {15};
\node[byte, right=of b0] (b1) {03};
\node[byte, right=of b1] (b2) {03};
\node[byte, right=of b2] (b3) {00};
\node[byte, right=of b3] (b4) {02};
\node[byte, right=of b4, fill=unknownbg] (b5) {02};
\node[byte, right=of b5, fill=unknownbg] (b6) {78};
\node[below=0.2cm of b0, font=\scriptsize, align=center] {ContentType\\alert (21)};
\node[font=\scriptsize, align=center] at ($(b1.south)!0.5!(b2.south)+(0,-0.55)$) {версия\\0x0303};
\node[font=\scriptsize, align=center] at ($(b3.south)!0.5!(b4.south)+(0,-0.55)$) {дължина\\2};
\node[below=0.2cm of b5, font=\scriptsize, align=center] {ниво\\fatal (2)};
\node[below=0.2cm of b6, font=\scriptsize, align=center] {описание\\120};
\end{tikzpicture}
\end{center}
\vspace{0.8em}
\begin{itemize}
\item \code{no\_application\_protocol} \src{RFC 8446, раздел 6; RFC 7301, раздел 3.2}
\item Интеграционният тест сравнява точно тези октети, а не текста на грешката. На Linux
  Python с OpenSSL 3.0.13 не може дори да назове предупреждението.
\end{itemize}
\end{frame}

% ---------------------------------------------------------------------------------
\begin{frame}{Заплахи и отговори}
\small
\renewcommand{\arraystretch}{1.2}
\begin{tabular}{@{}p{0.36\textwidth}p{0.36\textwidth}p{0.22\textwidth}@{}}
\toprule
\textbf{Заплаха} & \textbf{Отговор в кода} & \textbf{Проверено} \\
\midrule
свързва се и мълчи & общ краен срок от приемането & тест \\
праща \hex{0x16} и спира & същият срок покрива ръкостискането & тест \\
по един октет преди изтичане & срокът е абсолютен & по построение \\
много мълчащи връзки & граница; новите се затварят & тест \\
\hex{0x16} и безсмислени данни & OpenSSL отхвърля записа & тест \\
двоични данни, \code{get} & „непознато“, без отговор & тестове \\
огромна заглавна част & граница 8192 октета, отговор 431 & тест \\
стара версия на TLS & минимална версия на контекста & тест \\
\bottomrule
\end{tabular}\par\medskip
Разпознаването \textbf{насочва}, а не валидира. Проверките на записа ги прави OpenSSL.
\end{frame}

% ---------------------------------------------------------------------------------
\begin{frame}{Грешката, която тестовете откриха}
\begin{columns}[T]
\column{0.55\textwidth}
Първата версия проверяваше границата от 8192 октета само \emph{докато търси} празния ред,
който завършва заглавната част.\par\medskip
Едно четене донася и октетите над границата, и празния ред: празният ред се намира
пръв и заявката се приема.\par\medskip
\textbf{Поправка:} позицията на празния ред се сравнява с границата.
\column{0.4\textwidth}
\begin{tikzpicture}[node distance=0.35cm]
\node[box, text width=4.4cm] (u) {тест на функция: заявката по един октет\\\textcolor{tls}{\bfseries минава}};
\node[unkbox, text width=4.4cm, below=of u] (i) {интеграционен тест: поле от 9000 октета\\\textbf{200 вместо 431}};
\node[box, text width=4.4cm, below=of i] (n) {нов тест на функция: заявката на големи части};
\end{tikzpicture}
\end{columns}
\vspace{0.8em}
{\small Тестът на функция не можеше да я открие, защото не доставя данните така, както мрежата.\par}
\end{frame}

% ---------------------------------------------------------------------------------
\begin{frame}{Проверка}
\begin{columns}[T]
\column{0.5\textwidth}
\renewcommand{\arraystretch}{1.2}
\begin{tabular}{@{}lr@{}}
\toprule
\textbf{Тест} & \textbf{проверки} \\
\midrule
разпознаване, 256 стойности & \checksClassify \\
четене с връщане & \checksReplay \\
анализатор на HTTP & \checksHttp \\
TLS сесия в паметта & \checksTls \\
интеграционен, истински сокети & \checksIntegration \\
\bottomrule
\end{tabular}
\column{0.46\textwidth}
\begin{itemize}
\item Минават на \textbf{Windows 11} (MSVC 19.51, OpenSSL 3.6.4) и на \textbf{Ubuntu 24.04 във WSL2} (GCC 14.2, OpenSSL 3.0.13).
\item Без предупреждения при \code{/W4} и при \code{-Wall -Wextra -Wpedantic}.
\item macOS: конфигурация има, проверка няма.
\end{itemize}
\end{columns}
\end{frame}

% ---------------------------------------------------------------------------------
\begin{frame}{Измерване: постановка}
Четири варианта, всеки като отделен процес на една машина:\par\medskip
\begin{tabular}{@{}lll@{}}
\toprule
 & \textbf{с разпознаване} & \textbf{отделен слушател} \\
\midrule
\textcolor{tls}{\textbf{TLS}} & един сокет, решение по октета & ръкостискане веднага \\
\textcolor{clear}{\textbf{HTTP}} & един сокет, решение по октета & направо към анализатора \\
\bottomrule
\end{tabular}\par\medskip
\begin{itemize}
\item Една заявка на връзка. Клиент на Python, \winRate{} връзки в секунда, проверка на сертификата и името.
\item \winReps{} повторения по \winConns{} връзки на вариант, ред на вариантите разбъркан (начално число \winSeed{}).
\item Ефект = разлика на медианите; 95\% интервал с повторна извадка, 2000 пъти, в рамките на всяко повторение.
\item Минути, а не часове: Windows \winDuration{}~s, Linux \linDuration{}~s.
\end{itemize}
\end{frame}

% ---------------------------------------------------------------------------------
\begin{frame}{Първото измерване мереше друго}
\begin{columns}[T]
\column{0.55\textwidth}
Времената се мереха от \code{accept()}. Разликите излязоха
$\winfirstTlsDiffA$~µs за TLS и $\winfirstClearDiffA$~µs за HTTP.\par\medskip
Но при HTTP с отделен слушател времето до решението беше
\textbf{\winfirstClearDedicatedDecideMedian{}~ms}, а този вариант не чете нищо преди решението.\par\medskip
Мереше се стартирането на нишката, не разпознаването.
\column{0.4\textwidth}
\begin{tikzpicture}[node distance=0.3cm]
\node[box, text width=4.2cm] (a) {\code{accept()}};
\node[unkbox, text width=4.2cm, below=of a] (t) {стартиране на нишката\\\scriptsize еднакво за всички, шум};
\node[tlsbox, text width=4.2cm, below=of t] (s) {от тук се мери във второто измерване};
\draw[arr] (a) -- (t); \draw[arr] (t) -- (s);
\end{tikzpicture}
\end{columns}
\vspace{0.8em}
{\small Първото измерване е запазено в хранилището непроменено, като доказателство за проблема.\par}
\end{frame}

% ---------------------------------------------------------------------------------
\begin{frame}{Резултат}
\begin{columns}[T]
\column{0.58\textwidth}
{\small Ефект на разпознаването, от стартирането на нишката (µs, 95\% интервал)\par}\smallskip
\begin{tikzpicture}[x=0.085cm, y=0.62cm]
\def\xmin{-15}\def\xmax{60}
\draw[rule] (\xmin,0.4) -- (\xmax,0.4);
\foreach \t in {-10,0,10,20,30,40,50}{
  \draw[rule] (\t,0.4) -- (\t,-4.2);
  \node[font=\tiny, color=muted] at (\t,-4.5) {\t};
}
\draw[ink, line width=0.9pt] (0,0.4) -- (0,-4.2);
\effect{-0.5}{Windows, TLS}{\plotwinTlsLow}{\plotwinTlsHigh}{tls!45}
\fill[tls] (\plotwinTlsDiff,-0.5) circle (2.6pt);
\effect{-1.5}{Windows, HTTP}{\plotwinClearLow}{\plotwinClearHigh}{clear!45}
\fill[clear] (\plotwinClearDiff,-1.5) circle (2.6pt);
\effect{-2.5}{Linux, TLS}{\plotlinTlsLow}{\plotlinTlsHigh}{tls!45}
\fill[tls] (\plotlinTlsDiff,-2.5) circle (2.6pt);
\effect{-3.5}{Linux, HTTP}{\plotlinClearLow}{\plotlinClearHigh}{clear!45}
\fill[clear] (\plotlinClearDiff,-3.5) circle (2.6pt);
\end{tikzpicture}
\column{0.38\textwidth}
\textbf{В паметта}\par
Windows \microWinDiff{}~ns, Linux \microLinDiff{}~ns на връзка.\par\medskip
\textbf{TLS връзката}\par
\winTlsDedicatedMedian{}~ms на Windows, \linTlsDedicatedMedian{}~ms на Linux.\par\medskip
Цената е под една хилядна от връзката и под разделителната способност от край до край.
\end{columns}
\vspace{0.6em}
{\small Linux, HTTP: $\linClearDiffS$~µs с интервал без нулата. Причината не е установена
и числото не се приписва на разпознаването.\par}
\end{frame}

% ---------------------------------------------------------------------------------
\begin{frame}{Сериозното ограничение не е в кода}
\begin{columns}[T]
\column{0.55\textwidth}
Съчетаването на некриптиран и TLS трафик отваря пътя за \alert{SSL Stripping}: нападател
по средата премахва преминаването към TLS. \src{RFC 9325, 3.2; RFC 7457, 2.1}\par\medskip
Разпознаването не създава проблема (той съществува и при два порта), но прави
некриптирания достъп възможен по подразбиране.\par\medskip
\code{octet} сега: едно и също съдържание и по двата пътя, без HSTS.
\column{0.4\textwidth}
\textbf{Правилната схема}\par\smallskip
\begin{tikzpicture}[node distance=0.35cm]
\node[clearbox, text width=4.2cm] (c) {некриптиран път:\\само пренасочване към TLS};
\node[tlsbox, text width=4.2cm, below=of c] (t) {TLS път: съдържание\\и заглавно поле HSTS};
\draw[arr] (c) -- (t);
\end{tikzpicture}\par\smallskip
{\scriptsize HSTS: RFC 6797. Ползата от един порт остава.\par}
\end{columns}
\vspace{0.8em}
{\small Не е правено: fuzzing, одит от трета страна. Нишка на връзка не издържа на голяма атака за отказ от обслужване.\par}
\end{frame}

% ---------------------------------------------------------------------------------
\begin{frame}{Изводи}
\begin{itemize}
\item Един октет разделя TLS от HTTP за всичките 256 стойности. Основата е в стандартите.
\item Реализацията не губи данни и не отслабва TLS: връщане на октетите, общ краен срок, строг ALPN.
\item Цената е под една микросекунда на връзка. Производителността не е аргумент срещу механизма.
\item Аргументът, който остава, е сигурността: пренасочване и HSTS, иначе схемата е уязвима за SSL Stripping.
\end{itemize}
\vspace{0.8em}
\textbf{Следващи стъпки:} пренасочване и HSTS; fuzzing на анализатора и на разпознаването;
измерване на физически Linux през истинска мрежа.
\end{frame}

% ---------------------------------------------------------------------------------
\begin{frame}[plain]
\centering
\vfill
{\Huge\bfseries Въпроси?\par}
\vspace{1.5em}
{\large\ttfamily github.com/Alex-Tsvetanov/octet\par}
\vspace{0.6em}
{\small\color{muted} Кодът, тестовете, суровите записи от измерванията и документът\par}
\vfill
\end{frame}

% --- Резервни слайдове, за въпроси ------------------------------------------------
\appendix

\begin{frame}{Резерва: Windows, от стартирането на нишката (ms)}
\centering\small
% Generated by bench/analyse.py. Do not edit by hand.
% metric: served_from_start_us, milliseconds
\begin{tabular}{lrrrrr}
\toprule
Рамо & $n$ & p50 & p90 & p99 & медиани по повторения \\
\midrule
TLS, с разпознаване & 2000 & 3{,}137 & 6{,}964 & 20{,}840 & 3{,}088\,\dots\,3{,}236 \\
TLS, отделен слушател & 2000 & 3{,}108 & 6{,}254 & 22{,}496 & 3{,}076\,\dots\,3{,}154 \\
HTTP, с разпознаване & 2000 & 0{,}126 & 0{,}436 & 20{,}424 & 0{,}124\,\dots\,0{,}128 \\
HTTP, отделен слушател & 2000 & 0{,}128 & 0{,}502 & 18{,}616 & 0{,}117\,\dots\,0{,}136 \\
\bottomrule
\end{tabular}
\par\bigskip
{\footnotesize Последната колона: най-малката и най-голямата медиана сред повторенията.
При TLS разликата между тях е по-голяма от самите ефекти.\par}
\end{frame}

\begin{frame}{Резерва: Linux във WSL2, от стартирането на нишката (ms)}
\centering\small
% Generated by bench/analyse.py. Do not edit by hand.
% metric: served_from_start_us, milliseconds
\begin{tabular}{lrrrrr}
\toprule
Рамо & $n$ & p50 & p90 & p99 & медиани по повторения \\
\midrule
TLS, с разпознаване & 2000 & 2{,}169 & 13{,}636 & 27{,}076 & 2{,}084\,\dots\,4{,}008 \\
TLS, отделен слушател & 2000 & 2{,}163 & 10{,}390 & 16{,}126 & 2{,}109\,\dots\,9{,}543 \\
HTTP, с разпознаване & 2000 & 0{,}084 & 0{,}116 & 0{,}242 & 0{,}072\,\dots\,0{,}095 \\
HTTP, отделен слушател & 2000 & 0{,}079 & 0{,}110 & 0{,}183 & 0{,}068\,\dots\,0{,}090 \\
\bottomrule
\end{tabular}
\par\bigskip
{\footnotesize Ограничения: една машина, без мрежа; клиент на Python, по-бавен от сървъра;
Linux във виртуална машина; нишка на връзка.\par}
\end{frame}

\begin{frame}{Резерва: отклонение от RFC 9112}
RFC 9112, раздел 2.2: сървърът \emph{SHOULD} да игнорира празен ред преди реда на заявката.\par\medskip
\code{octet} не го прави: връзка, която започва с \hex{CR (0x0D)}, се затваря като непозната.\par\medskip
Причина: иначе решението не може да се вземе по първия октет и трябва да се чака края на
празните редове. Изборът е съзнателен и е описан в документа.
\end{frame}

\end{document}
